\section{Introduction}

Adding one generator and one relation to a nontrivial group should leave
a nontrivial group. This is the Kervaire conjecture. We prove it by
establishing coefficient injectivity in the case where the new generator
has exponent sum $\pm1$.

Let $A$ be a group, let $\langle t\rangle\cong\Z$, and set
$H=A*\langle t\rangle$. Write $p:H\to\Z$ for the homomorphism that kills
$A$ and sends $t$ to $1$. A word $w\in H$ is \emph{unimodular} if
$p(w)=\pm1$. Denote its normal closure in $H$ by
$\langle\!\langle w\rangle\!\rangle_H$.

The equation $w(t)=1$ is \emph{solvable over $A$} if there are an
embedding $A\to B$ and an element $b\in B$ such that evaluating $t$ at
$b$ gives $w(b)=1$. This is equivalent to injectivity of the coefficient
map $A\to H/\langle\!\langle w\rangle\!\rangle_H$: every solution
factors through this quotient, while an injective coefficient map lets
the quotient itself supply a solution.

\Needspace{7\baselineskip}
\begin{samepage}
\begin{theorem}\label{thm:main}
For every group $A$ and every unimodular word $w\in A*\langle t\rangle$,
the canonical homomorphism
\[
 A\longrightarrow (A*\langle t\rangle)/\langle\!\langle w\rangle\!\rangle_H
\]
is injective.
\end{theorem}
\end{samepage}

No restriction on torsion, cardinality, or presentation of $A$ is imposed.
The nontriviality assertion for all exponent sums follows at once.

\begin{corollary}[The Kervaire conjecture]\label{cor:kervaire}
If $A\ne1$, then
$(A*\langle t\rangle)/\langle\!\langle w\rangle\!\rangle_H\ne1$
for every $w\in A*\langle t\rangle$.
\end{corollary}
\begin{proof}
For $p(w)=\pm1$, apply Theorem~\ref{thm:main}.
For every other integer $d=p(w)$, the map $p$ induces a surjection from
the quotient onto the nontrivial group $\Z/d\Z$, where $\Z/0\Z=\Z$.
\end{proof}

\subsection{Context and methods}

Kervaire's characterization of high-dimensional knot groups includes
normal generation by a single element \cite[Theorem 1.1]{Kervaire1965}.
Chen \cite[preprint Section 1]{Chen2026} traces the group-theoretic
Kervaire problem to this setting. Its stronger Kervaire--Laudenbach form
asks for coefficient injectivity whenever $p(w)\ne0$; such one-variable
equations are called \emph{nonsingular}. Theorem~\ref{thm:main} proves its
unimodular case, while Corollary~\ref{cor:kervaire} resolves the full
nontriviality conjecture.
At the level of assertions for all groups, this nontriviality statement
is equivalent to unimodular coefficient injectivity: Klyachko records the
standard reduction through an embedding into a simple group
\cite[preprint Section 3, Proposition 2]{Klyachko2004}.
We prove injectivity directly.

Gerstenhaber and Rothaus proved solvability of nonsingular systems of
equations over finite groups in finite overgroups
\cite[Theorem 2]{GerstenhaberRothaus1962}.
In particular their theorem gives the required injectivity for finite
coefficient groups, and hence for residually finite groups by separating
each nonidentity coefficient in a finite quotient.
Pestov records the resulting ultraproduct extension: nonsingular
equations over hyperlinear groups have solutions in hyperlinear
overgroups \cite[preprint Corollary 10.4]{Pestov2008}.
Hyperlinear groups admit embeddings into metric ultraproducts of unitary
groups with normalized Hilbert--Schmidt metrics.
Klyachko proved unimodular injectivity for torsion-free groups
\cite[Theorem 4, exponent-one case]{Klyachko1993}. Chen gave a new proof
through estimates for the complexity of surfaces
\cite[preprint Theorem 6.9]{Chen2026}. For arbitrary coefficient groups,
Klyachko and Lurye proved injectivity after imposing $w^m=1$ when $w$ is
unimodular and $m\geq2$
\cite[preprint Section 1, Theorem; Section 5]{KlyachkoLurye2012}; this does not impose
$w=1$.

Kawauchi has published a proposed resolution of the universal
nontriviality assertion through knot exteriors, relative collapses, and
ribbon sphere-link groups \cite[Theorem 1 and Corollary 2]{Kawauchi2024}.
By the equivalence just recalled, its endpoint includes unimodular
injectivity. The argument below takes a different route through spectral
phase and a finite-dimensional fixed-space theorem; neither the collapse
construction nor the ribbon-asphericity assertions enter its proof.
Recent work of Klyachko, Mikheenko, and Olshanskii also distinguishes
solvability inside a coefficient group from solvability in an overgroup
\cite[Sections 1 and 3]{KlyachkoMikheenkoOlshanskii2026}.
Our conclusion concerns the latter and does not assert that a solution
already lies in $A$.

Unitary groups and degree theory already play a central role in the
Gerstenhaber--Rothaus method \cite[Theorem 1]{GerstenhaberRothaus1962}; see also
\cite[Section 1.1]{KlyachkoThom2017}.
Corner-labelled diagrams and Euler-characteristic estimates likewise
have a substantial history in equations over groups
\cite[preprint Sections 3--4]{KlyachkoLurye2012}
and \cite[preprint Sections 2--4]{Chen2026}.
Here these methods meet through a trace comparison for permutations
weighted by the left regular representation of $A$. The regular
representation is available for every group and supplies a faithful
finite trace without an approximation hypothesis on $A$.
For the spectral ingredient, Thompson's matrix exponential formula and
its extensions provide a related precedent
\cite[Section 2 and Theorem 4.4]{AntezanaLarotondaVarela2012};
Section~\ref{sec:phase} distinguishes those formulas from the unrestricted
phase inequality proved here.

Two ingredients are useful independently of the diagram argument.
The first is subadditivity of the trace of the fractional spectral phase
of a unitary. The second is a fixed-space lemma: given $k\geq1$, signs
$s_1,\ldots,s_n\in\{1,-1\}$ with $\sum_i s_i=\pm1$, and $P\in\U(nk)$,
some $X\in\U(k)$ makes $\diag(X^{s_1},\ldots,X^{s_n})P$
have at least $k$ fixed directions. We prove the latter by a degree
calculation on a compact manifold that records a unitary together with
a fixed $k$-plane. Their combination gives the planar obstruction below.

\subsection{Proof overview}

Suppose a coefficient $g\in A$ becomes trivial after imposing $w=1$.
A finite expression for $g$ as a product of conjugates of $w$ and $w^{-1}$
can be drawn on a punctured disk. Pairing the occurrences of $t$ produces
disjoint bands between word disks. The complementary boundary curves
read words in $A$.

The central assertion is Theorem~\ref{thm:planar}: in a connected planar
surface made from these disks and bands, one boundary word cannot be the
only nontrivial one. To prove it, we encode travel along the boundaries
by a weighted permutation. We measure that operator by the trace of its
spectral phase. Mixing positive and negative word disks with a suitable
unitary gives an upper bound. A second calculation, one boundary circle
at a time, approaches that bound plus the nonnegative phase of the
exceptional word. The Euler characteristic formula identifies the common
term in these two calculations. The upper bound forces the extra term to
vanish, so this word is trivial.

Section~\ref{sec:phase} proves the phase inequality, and
Section~\ref{sec:degree} proves the fixed-space lemma used to choose the
mixing unitary. Section~\ref{sec:diagram} establishes the planar boundary
theorem. Section~\ref{sec:topology} constructs the bands from the original
group relation and removes their connected components from the inside
out, proving $g=1$.

The proof uses group presentations, the spectral theorem for unitary
operators, elementary surface topology, and degree modulo two. The trace,
incidence-space, and diagram calculations are included.
